\documentclass[10pt,a4paper]{amsart}
\usepackage[margin=19mm]{geometry}
\usepackage{amsmath,amssymb,mathtools}
\usepackage[scr=rsfs,cal=euler]{mathalfa}
\usepackage{xfrac}
\usepackage[unicode,hidelinks,bookmarks=false]{hyperref}
\newcommand{\ie}{i.e.\ }
\newcommand{\cf}{cf.\ }
\newcommand{\resp}{respectively\ }
\newcommand{\Subsection}{\subsection}
\providecommand{\nobreakdash}{\nobreak\hbox{-}}
\setlength{\emergencystretch}{2em}
\multlinegap0pt
\usepackage[all]{xy}
\usepackage{amssymb}
\usepackage{enumerate}
\newtheorem{lem}{Lemma}
\title{Tensor Abelian Geometry of VI-modules}
\author{Peng Xu}
\date{Source: arXiv:2601.14020v2 (CC BY 4.0)}
\begin{document}
\maketitle

\noindent\emph{Source and licence.} Tensor Abelian Geometry of VI-modules. Version arXiv:2601.14020v2 (CC BY 4.0), distributed by arXiv under the Creative Commons Attribution 4.0 International licence, http://creativecommons.org/licenses/by/4.0/, as recorded in the arXiv licence metadata for this exact version and shown on the abstract page https://arxiv.org/abs/2601.14020v2. Full licence notice: \texttt{licenses/arxiv-2601.14020-CC-BY-4.0.txt}.\par\smallskip

\noindent\textit{Editorial scope.} Counted unit: the article's lem lem4.13 (source0.tex lines 475--476). The article's complete original proof of it is delivered with it (source0.tex lines 477--496, one \texttt{proof} environment of 20 lines). No mathematics is written, repaired or re-derived: the statement and the proof are the article's own lines, in the article's own notation, and no retained line is substituted or re-flowed.  Retained uncounted as the source-local context the proof needs: lines 461--470.  Nothing was omitted from any retained span, so the package carries no omission record.  No retained line contains a citation, so the wrapper carries no bibliography and no bibliography entry.  No retained line contains a figure environment, an image-inclusion command or drawing code, so no image is delivered and the package declares no asset dependency.  Editorial formatting only, in addition: an AMS \texttt{amsart} preamble that restates the article's own declarations and macros used by the retained lines, namely the environments \texttt{lem} on separate counters, together with the standard packages those lines need.  The retained environments print in this document as Lemma 1 in order of appearance, whereas the article prints the same items with the numbering of its own full document.  The complete native source of the article as distributed in the arXiv e-print for this exact version is bundled unchanged as \texttt{sources/193063/source0.tex}.
\begin{lem}\label{lem4.12} Let $i:\mathcal{U}\rightarrow \mathcal{V}$ be the inclusion of a down-closed subcategory  and $j:\mathcal{V}\setminus~\mathcal{U}\rightarrow \mathcal{V}$ be the complement. Suppose $\mathcal{V}$ is noetherian, we have:
\begin{enumerate}
\item $j_{!}j^{*}X\rightarrow X\rightarrow i_{*}i^{*}X$ is short exact for any $X\in \mathrm{A}(\mathcal{V})^{c}$ and thus $\mathrm{A}(\mathcal{V})^{c}/\mathrm{A}(\mathcal{V}\setminus \mathcal{U})^{c}\simeq_{\otimes} \mathrm{A}(\mathcal{U})^{c} $;

\item if $X\in \mathrm{A}(\mathcal{V})^{c}$ and $ \mathrm{supp}(X)\subseteq \mathcal{V}\setminus~\mathcal{U}$, then $j_{!}j^{*}X\cong X$;

\item $X\in \mathrm{Serre}_{\otimes}\langle e_{G}|~G\in \uparrow(\mathrm{supp}(X))\rangle$ for any $X\in \mathrm{A}(\mathcal{V})^{c}$.

\end{enumerate}
\end{lem}

\begin{lem}\label{lem4.13} Let $i:\mathcal{U}\rightarrow \mathcal{V}$ be the inclusion of a down-closed subcategory  and $j:\mathcal{V}\setminus~\mathcal{U}\rightarrow \mathcal{V}$ be the complement. Suppose $\mathcal{V}$ is noetherian and $X\in \mathrm{A}(\mathcal{U})^{c}$, then we have: $$i_{!}\mathrm{Serre}_{\otimes}\langle X\rangle\subseteq \mathrm{Serre}_{\otimes}\langle \{i_{!}X\}\cup \{e_{G}|~G\in \mathcal{V}\setminus~\mathcal{U}\}\rangle.$$
\end{lem}

\begin{proof} Firstly we claim that $i_{!}(X\otimes Y)\in \mathrm{Serre}_{\otimes}\langle \{i_{!}X\}\cup \{e_{G}|~G\in \mathcal{V}\setminus~\mathcal{U}\}\rangle$ for any $Y\in \mathrm{A}(\mathcal{U})^{c}$.  Consider the following diagram: $$\xymatrix{
                & i_{!}(X\otimes Y) \ar@{->>}[dr] \ar[rr]^{h} & &i_{!}(X)\otimes i_{!}(Y) \ar@{->>}[rr] & & \mathrm{Coker}(h)         \\
 \mathrm{Ker}(h) \ar@{>->}[ur] & &     \mathrm{Im}(h)\ar@{>->}[ur]       }
$$
where the map $h$ is induced by the oplax monoidal structure on $i_{!}$. Since $i^{*}$ is strong monoidal and exact, we have $i^{*}\mathrm{Coker}(h)=0$ and so $\mathrm{Coker}(h)\in \mathrm{A}(\mathcal{U})^{c}$ by Lemma \ref{lem4.12}. Thus $$\mathrm{Coker}(h)\cong j_{!}j^{*}\mathrm{Coker}(h)\in \mathrm{Serre}_{\otimes}\langle \{e_{G}|~G\in \mathcal{V}\setminus~\mathcal{U}\}\rangle$$ by the same lemma again. Similarly $i^{*}\mathrm{Ker}(h)=0$ and so $$\mathrm{Ker}(h)\cong j_{!}j^{*}\mathrm{Ker}(h)\in \mathrm{Serre}_{\otimes}\langle \{e_{G}|~G\in \mathcal{V}\setminus~\mathcal{U}\}\rangle.$$
Thus \begin{align*}
  i_{!}(X\otimes Y) &\in \mathrm{Serre}_{\otimes}\langle \mathrm{Ker}(h), \mathrm{Im}(h), \mathrm{Coker}(h) \rangle\\
  &\subseteq \mathrm{Serre}_{\otimes}\langle \mathrm{Ker}(h), i_{!}(X)\otimes i_{!}(Y) \rangle\\
  &\subseteq \mathrm{Serre}_{\otimes}\langle  \{i_{!}X\}\cup \{e_{G}|~G\in \mathcal{V}\setminus~\mathcal{U}\} \rangle,
\end{align*}
as claimed.

Now suppose $Y\rightarrowtail X\twoheadrightarrow Z$ is short exact. Then $$ i_{!}(Z)\in \mathrm{Serre}_{\otimes}\langle  i_{!}X \rangle\subseteq \mathrm{Serre}_{\otimes}\langle  \{i_{!}X\}\cup \{e_{G}|~G\in \mathcal{V}\setminus~\mathcal{U}\} \rangle$$ since $i_{!}$ is right exact. Apply $i_{!}$ to the injection $f:Y\rightarrowtail X$, we have $i^{*}\mathrm{Ker}(i_{!}f)=0$ and so $\mathrm{Ker}(i_{!}f)\in \mathrm{Serre}_{\otimes}\langle \{e_{G}|~G\in \mathcal{V}\setminus~\mathcal{U}\rangle$ as above. Then $$i_{!}(Y)\in \mathrm{Serre}_{\otimes}\langle \mathrm{Ker}(i_{!}f), \mathrm{Im}(i_{!}f)\rangle \subseteq \mathrm{Serre}_{\otimes}\langle  \{i_{!}X\}\cup \{e_{G}|~G\in \mathcal{V}\setminus~\mathcal{U}\} \rangle.$$

At last, for any short exact sequence $Y\rightarrowtail W\twoheadrightarrow Z$ in $\mathrm{A}(\mathcal{U})^{c}$ with $i_{!}(Y), i_{!}(Z)\in  \mathrm{Serre}_{\otimes}\langle  \{i_{!}X\}\cup \{e_{G}|~G\in \mathcal{V}\setminus~\mathcal{U}\} \rangle$, then \begin{align*}
  i_{!}(W) &\in \mathrm{Serre}_{\otimes}\langle  \mathrm{Im}(g), i_{!}(Z) \rangle\\
  &\subseteq \mathrm{Serre}_{\otimes}\langle i_{!}(Y), i_{!}(Z) \rangle\\
  &\subseteq \mathrm{Serre}_{\otimes}\langle  \{i_{!}X\}\cup \{e_{G}|~G\in \mathcal{V}\setminus~\mathcal{U}\} \rangle
\end{align*} where $g:Y\rightarrow W$ is the monomorphism. This concludes the proof.
\end{proof}

\end{document}
